% !TeX encoding = UTF-8
% !TeX program = xelatex
\documentclass{resume}
\usepackage{zh_CN-Adobefonts_external}
\usepackage{linespacing_fix}
\usepackage[English]{languageSelection}
\titlespacing*{\section}{0cm}{*0.65}{*0.45}
\titlespacing*{\subsection}{0cm}{*0.45}{*0.2}

\begin{document}
\pagenumbering{gobble}
\fontsize{10.2pt}{11.8pt}\selectfont

\vspace{-10pt}
\name{Yuhan Liu}

\info{+86 18771182913~|~yuhanliu050128@gmail.com}{}{}{}
\info{https://hefengchuyang-sketch.github.io/}{}{}{}
\info{GitHub: github.com/hefengchuyang-sketch/POUW-Chain}{}{}{}

\vspace{2pt}
\centerline{\small Research Interests: Agentic AI | Distributed AI Systems | Trustworthy Computing}

\vspace{1pt}
\section{Publications}
\datedsubsection{\textbf{FVEN: Fluid Velocity Estimation Network Based on Laser-Induced Plasma Images}}{2025}
\textbf{First Author}

Asian Conference on Artificial Intelligence Technology (EI Indexed)

\vspace{3pt}
\datedsubsection{\textbf{Fluid Velocity Measurement using Laser-induced Plasma Tracking Velocimetry and Deep Learning}}{2026}
\textbf{Co-Author}

Measurement (SCI Q2, IF=5.0)

\section{Education}
\datedsubsection{\textbf{Taiyuan University of Technology}}{2023.09--2027.06}
B.Eng. in Artificial Intelligence

Expected Graduation: Jun 2027

Yunding Academy (AI Track Member)

Mihrilab (Project Lead)

\vspace{1pt}
\section{Honors and Awards}
\datedsubsection{27th China Robot and Artificial Intelligence Competition ~ Third Prize}{2025}
\vspace{0.5pt}
\datedsubsection{5th Wutong Cup National College AI Innovation Competition, Agent Track ~ National First Prize}{2026}

\vspace{1pt}
\section{Research Experience}
\noindent{\large \textbf{\parbox[t]{0.78\textwidth}{FVEN: Fluid Velocity Estimation Network Based on Laser-Induced Plasma Images (First Author)}}\hfill 2024.06 -- 2025.12\par}
\textbf{Contributions}

\begin{itemize}
\item Proposed a three-stage FVEN framework.
\item Designed CEM and MIAAM modules.
\item Built datasets containing 7,200 airflow and 1,440 combustion images.
\item Achieved MAE $\leq$ 3.17 m/s and relative error $\leq$ 4\%.
\item Published as first author in ACAIT (EI).
\end{itemize}

\vspace{1pt}
\vspace{0.8\baselineskip}
\section{Open-Source Project}
\datedsubsection{\textbf{POUW-Chain: Decentralized AI Computation Platform}}{2025.12 -- 2026.04}
\textbf{GitHub:} https://github.com/hefengchuyang-sketch/POUW-Chain

\textbf{Positioning:} Designed and implemented a decentralized AI computation platform for trustworthy execution, resource scheduling, and verifiable settlement.

\textbf{Key Contributions}

\begin{itemize}
\item Designed useful-work consensus mechanisms.
\item Built heterogeneous compute scheduling.
\item Developed verifiable execution and settlement protocols.
\item Implemented 60+ modules and validated functionality through 298 automated tests.
\end{itemize}

\vspace{2pt}
\vspace{0.8\baselineskip}
\section{Projects}
\datedsubsection{\textbf{Competition Project: Intelligent Site Selection Master}}{2026.02}
\textbf{Role:} Team member and lead agent developer

\textbf{Award:} National First Prize, Agent Track, 5th Wutong Cup National College AI Innovation Competition

\textbf{Research Contributions}

\begin{itemize}
\item Built a five-agent collaborative system powered by Qwen2.5-7B and Agentic-RAG.
\item Developed Agentic-RAG pipelines.
\item Integrated geospatial and mobility datasets.
\item Achieved 89\% Top-3 recommendation accuracy.
\end{itemize}

\vspace{1pt}
\section{Skills}
Python, Solidity, PyTorch, FastAPI, Agentic AI, Multi-Agent Systems, Distributed Systems, Blockchain, Trusted Execution Environments (TEE)
\end{document}
